On PhysioNet 2012 set-a, with similar-sized tuning grids (\rhval{const/tune_cells_lr}, \rhval{const/tune_cells_gbdt} and \rhval{const/tune_cells_gru} cells for LR, GBDT and each GRU) and \rhval{const/n_seeds} seeds of \rhval{const/n_folds}-fold patient-level cross-validation, gradient boosting on per-variable summaries (AUROC \rhval{main/gbdt-summary/physionet2012_mortality/auroc/mean:3}) beat logistic regression (\rhval{main/lr-summary/physionet2012_mortality/auroc/mean:3}) and all three recurrent models (\rhval{main/gru-d/physionet2012_mortality/auroc/mean:3} for GRU-D), and GRU-D did not beat a forward-filled GRU or a GRU with mask and delta inputs at the full 48 h horizon. The decay mechanism of GRU-D had no measurable effect in our ablations. Missingness-aware inputs did help the GRU at shorter horizons (\rhval{const/h12} h and \rhval{const/h24} h), and explicit count features helped logistic regression but not boosting. These conclusions are limited to one small cohort, our reimplementations, hourly binning and a modest tuning budget; they do not show that missingness-aware recurrent models are worse in general, only that on this benchmark they were not better than a tuned tree ensemble on summaries.
